% GENERATED FILE. Edit canonical lessons, then run npm run generate:books.
% canonical-source-hash: fnv1a64:a987124b85a5017e
% canonical-lessons: KA-C434-vilamba, KA-C434-agamisu, KA-C434-nirgamisu, KA-C434-nilugade, KA-C434-vicarisu

\chapter{A Late Train}
\label{ch:ka-a-late-train}
\hlchaptermodality{434}

\begin{chapteropening}
\textbf{By the end of this chapter:} \emph{I can say a delay, to arrive, to depart, a halt and to enquire.}

Complete the last lesson of chapter 434: I can say a delay, to arrive, to depart, a halt and to enquire.
\end{chapteropening}

\section[\texorpdfstring{\kn{ವಿಳಂಬ} (viḷamba)}{ವಿಳಂಬ (viḷamba)}]{\kn{ವಿಳಂಬ} (viḷamba) — a delay}

\label{lesson:KA-C434-vilamba}



\begin{quote}
Before the new one: say the Kannada for a rail track, then the Kannada for to reach.
\end{quote}

\subsection*{You'll want to know: \kn{ವಿಳಂಬ}}
\textbf{\kn{ವಿಳಂಬ}} — \emph{viḷamba} — \textquotedblleft{}a delay\textquotedblright{}.

\begin{quote}
\kn{ಮಳೆಯ} \kn{ಕಾರಣ} \kn{ರೈಲು} \kn{ಎರಡು} \kn{ಗಂಟೆ} \kn{ವಿಳಂಬವಾಯಿತು}.
\end{quote}

\begin{quote}
\emph{maḷeya kāraṇa railu eraḍu gaṇṭe viḷambavāyitu.}
\end{quote}

\begin{quote}
Because of the rain the train was two hours late.
\end{quote}

With \kn{ಆಗು} it says that something ran late: \kn{ವಿಳಂಬವಾಯಿತು}, it was delayed. In speech people say \kn{ತಡವಾಯಿತು}.

\begin{cousinweb}
\textbf{\kn{ವಿಳಂಬ}} (\emph{viḷamba}) is borrowed from Sanskrit \emph{vilamba}, \textquotedblleft{}delay\textquotedblright{} (CDIAL 11890), from the root \emph{lamb}, \textquotedblleft{}to hang down, to lag behind\textquotedblright{}. Kannada writes its \emph{l} as \kn{ಳ}.
\end{cousinweb}

\begin{grammarlens}[title={What you've built}]
One more word for getting around.
\end{grammarlens}

\subsection*{Your turn}
\begin{itemize}
\raggedright
  \item \emph{Say it:} \emph{viḷamba}
  \item \emph{Say it:} \emph{viḷamba}, once more
  \item \emph{Say it:} \emph{maḷeya kāraṇa railu eraḍu gaṇṭe viḷambavāyitu.} — Because of the rain the train was two hours late.
  \item \emph{Recall:} say the Kannada for a rail track, then the Kannada for to reach, then say \emph{viḷamba} again
\end{itemize}

\begin{tcolorbox}[breakable,colback=teal!4,colframe=teal!35!black,title={Before you move on}]
What does \kn{ವಿಳಂಬ} mean? (\textquotedblleft{}A delay\textquotedblright{}.) Say it once more.
\end{tcolorbox}
\section[\texorpdfstring{\kn{ಆಗಮಿಸು} (āgamisu)}{ಆಗಮಿಸು (āgamisu)}]{\kn{ಆಗಮಿಸು} (āgamisu) — to arrive}

\label{lesson:KA-C434-agamisu}



\begin{quote}
Before the new one: say the Kannada for to reach, then the Kannada for a delay.
\end{quote}

\subsection*{You'll want to know: \kn{ಆಗಮಿಸು}}
\textbf{\kn{ಆಗಮಿಸು}} — \emph{āgamisu} — \textquotedblleft{}to arrive\textquotedblright{}.

\begin{quote}
\kn{ವಿಮಾನ} \kn{ಸಮಯಕ್ಕೆ} \kn{ಸರಿಯಾಗಿ} \kn{ಆಗಮಿಸಿತು}.
\end{quote}

\begin{quote}
\emph{vimāna samayakke sariyāgi āgamisitu.}
\end{quote}

\begin{quote}
The plane arrived right on time.
\end{quote}

Past (I …): \kn{ಆಗಮಿಸಿದೆ}, \emph{āgamiside}.

A formal verb of notices and announcements; in speech people say \kn{ಬಂದಿತು}, it came.

\begin{cousinweb}
\textbf{\kn{ಆಗಮಿಸು}} (\emph{āgamisu}) is built from Sanskrit \emph{āgama}, \textquotedblleft{}coming, arrival\textquotedblright{} (Monier-Williams), from \emph{ā-} \textquotedblleft{}towards\textquotedblright{} and \emph{gam} \textquotedblleft{}to go\textquotedblright{}, with the Kannada ending -\emph{isu} that makes a verb of a Sanskrit noun.
\end{cousinweb}

\begin{grammarlens}[title={What you've built}]
One more word for getting around.
\end{grammarlens}

\subsection*{Your turn}
\begin{itemize}
\raggedright
  \item \emph{Say it:} \emph{āgamisu}
  \item \emph{Say it:} \emph{āgamisu}, once more
  \item \emph{Say it:} \emph{vimāna samayakke sariyāgi āgamisitu.} — The plane arrived right on time.
  \item \emph{Recall:} say the Kannada for to reach, then the Kannada for a delay, then say \emph{āgamisu} again
\end{itemize}

\begin{tcolorbox}[breakable,colback=teal!4,colframe=teal!35!black,title={Before you move on}]
What does \kn{ಆಗಮಿಸು} mean? (\textquotedblleft{}To arrive\textquotedblright{}.) Say it once more.
\end{tcolorbox}
\section[\texorpdfstring{\kn{ನಿರ್ಗಮಿಸು} (nirgamisu)}{ನಿರ್ಗಮಿಸು (nirgamisu)}]{\kn{ನಿರ್ಗಮಿಸು} (nirgamisu) — to depart; to leave}

\label{lesson:KA-C434-nirgamisu}



\begin{quote}
Before the new one: say the Kannada for a delay, then the Kannada for to arrive.
\end{quote}

\subsection*{You'll want to know: \kn{ನಿರ್ಗಮಿಸು}}
\textbf{\kn{ನಿರ್ಗಮಿಸು}} — \emph{nirgamisu} — \textquotedblleft{}to depart; to leave\textquotedblright{}.

\begin{quote}
\kn{ರೈಲು} \kn{ಬೆಳಗ್ಗೆ} \kn{ಆರು} \kn{ಗಂಟೆಗೆ} \kn{ನಿರ್ಗಮಿಸುತ್ತದೆ}.
\end{quote}

\begin{quote}
\emph{railu beḷagge āru gaṇṭege nirgamisuttade.}
\end{quote}

\begin{quote}
The train departs at six in the morning.
\end{quote}

Past (I …): \kn{ನಿರ್ಗಮಿಸಿದೆ}, \emph{nirgamiside}.

Its partner is \kn{ಆಗಮಿಸು}, to arrive; station boards pair the two.

\begin{cousinweb}
\textbf{\kn{ನಿರ್ಗಮಿಸು}} (\emph{nirgamisu}) is built from Sanskrit \emph{nirgama}, \textquotedblleft{}departure, a way out\textquotedblright{} (CDIAL 7294), from \emph{nis-} \textquotedblleft{}out\textquotedblright{} and \emph{gam} \textquotedblleft{}to go\textquotedblright{}, with the verb ending -\emph{isu}.
\end{cousinweb}

\begin{grammarlens}[title={What you've built}]
One more word for getting around.
\end{grammarlens}

\subsection*{Your turn}
\begin{itemize}
\raggedright
  \item \emph{Say it:} \emph{nirgamisu}
  \item \emph{Say it:} \emph{nirgamisu}, once more
  \item \emph{Say it:} \emph{railu beḷagge āru gaṇṭege nirgamisuttade.} — The train departs at six in the morning.
  \item \emph{Recall:} say the Kannada for a delay, then the Kannada for to arrive, then say \emph{nirgamisu} again
\end{itemize}

\begin{tcolorbox}[breakable,colback=teal!4,colframe=teal!35!black,title={Before you move on}]
What does \kn{ನಿರ್ಗಮಿಸು} mean? (\textquotedblleft{}To depart; to leave\textquotedblright{}.) Say it once more.
\end{tcolorbox}
\section[\texorpdfstring{\kn{ನಿಲುಗಡೆ} (nilugaḍe)}{ನಿಲುಗಡೆ (nilugaḍe)}]{\kn{ನಿಲುಗಡೆ} (nilugaḍe) — a halt; a stop}

\label{lesson:KA-C434-nilugade}



\begin{quote}
Before the new one: say the Kannada for to arrive, then the Kannada for to depart.
\end{quote}

\subsection*{You'll want to know: \kn{ನಿಲುಗಡೆ}}
\textbf{\kn{ನಿಲುಗಡೆ}} — \emph{nilugaḍe} — \textquotedblleft{}a halt; a stop\textquotedblright{}.

\begin{quote}
\kn{ಈ} \kn{ರೈಲಿಗೆ} \kn{ಆ} \kn{ಊರಿನಲ್ಲಿ} \kn{ನಿಲುಗಡೆ} \kn{ಇಲ್ಲ}.
\end{quote}

\begin{quote}
\emph{ī railige ā ūrinalli nilugaḍe illa.}
\end{quote}

\begin{quote}
This train has no stop in that town.
\end{quote}

Built on \kn{ನಿಲ್ಲು}, to stop: a bus stop or a parking place can also be a \kn{ನಿಲುಗಡೆ}.

\begin{cousinweb}
\textbf{\kn{ನಿಲುಗಡೆ}} (\emph{nilugaḍe}) is a native Kannada noun built on \emph{nil}, \emph{nilu}, \textquotedblleft{}to stand, to stop\textquotedblright{} (DEDR 3675). Where the ending -\emph{gaḍe} comes from is not settled.
\end{cousinweb}

\begin{grammarlens}[title={What you've built}]
One more word for getting around.
\end{grammarlens}

\subsection*{Your turn}
\begin{itemize}
\raggedright
  \item \emph{Say it:} \emph{nilugaḍe}
  \item \emph{Say it:} \emph{nilugaḍe}, once more
  \item \emph{Say it:} \emph{ī railige ā ūrinalli nilugaḍe illa.} — This train has no stop in that town.
  \item \emph{Recall:} say the Kannada for to arrive, then the Kannada for to depart, then say \emph{nilugaḍe} again
\end{itemize}

\begin{tcolorbox}[breakable,colback=teal!4,colframe=teal!35!black,title={Before you move on}]
What does \kn{ನಿಲುಗಡೆ} mean? (\textquotedblleft{}A halt; a stop\textquotedblright{}.) Say it once more.
\end{tcolorbox}
\section[\texorpdfstring{\kn{ವಿಚಾರಿಸು} (vicārisu)}{ವಿಚಾರಿಸು (vicārisu)}]{\kn{ವಿಚಾರಿಸು} (vicārisu) — to enquire; to ask after}

\label{lesson:KA-C434-vicarisu}



\begin{quote}
Before the new one: say the Kannada for to depart, then the Kannada for a halt.
\end{quote}

\subsection*{You'll want to know: \kn{ವಿಚಾರಿಸು}}
\textbf{\kn{ವಿಚಾರಿಸು}} — \emph{vicārisu} — \textquotedblleft{}to enquire; to ask after\textquotedblright{}.

\begin{quote}
\kn{ನಾನು} \kn{ನಿಲ್ದಾಣದಲ್ಲಿ} \kn{ರೈಲಿನ} \kn{ಸಮಯ} \kn{ವಿಚಾರಿಸಿದೆ}.
\end{quote}

\begin{quote}
\emph{nānu nildāṇadalli railina samaya vicāriside.}
\end{quote}

\begin{quote}
I enquired about the train times at the station.
\end{quote}

Past (I …): \kn{ವಿಚಾರಿಸಿದೆ}, \emph{vicāriside}.

What is asked about takes -\kn{ಅನ್ನು} or no ending; \kn{ಆರೋಗ್ಯ} \kn{ವಿಚಾರಿಸು}, to ask after someone's health.

\begin{cousinweb}
\textbf{\kn{ವಿಚಾರಿಸು}} (\emph{vicārisu}) is built from Sanskrit \emph{vicāra}, \textquotedblleft{}deliberation, examination, investigation\textquotedblright{} (Monier-Williams), with the verb ending -\emph{isu}. The noun \kn{ವಿಚಾರ}, an idea or a matter, comes from the same Sanskrit word.
\end{cousinweb}

\begin{grammarlens}[title={What you've built}]
One more word for getting around.
\end{grammarlens}

\subsection*{Your turn}
\begin{itemize}
\raggedright
  \item \emph{Say it:} \emph{vicārisu}
  \item \emph{Say it:} \emph{vicārisu}, once more
  \item \emph{Say it:} \emph{nānu nildāṇadalli railina samaya vicāriside.} — I enquired about the train times at the station.
  \item \emph{Recall:} say the Kannada for to depart, then the Kannada for a halt, then say \emph{vicārisu} again
\end{itemize}

\begin{tcolorbox}[breakable,colback=teal!4,colframe=teal!35!black,title={Before you move on}]
What does \kn{ವಿಚಾರಿಸು} mean? (\textquotedblleft{}To enquire; to ask after\textquotedblright{}.) Say it once more.
\end{tcolorbox}
